\abstractch

乳腺癌已成为全球女性发病率最高的恶性肿瘤，早期准确诊断对提高患者生存率具有决定性意义。乳腺超声检查因其无创无辐射、实时成像、适合致密型乳腺等优势，已被广泛用作一线筛查手段，但传统的人工判读方式高度依赖医生经验，存在主观性强、一致性差、效率受限等问题。近年来，深度学习技术在医学影像分析领域取得了突破性进展，为开发自动化的乳腺癌超声图像计算机辅助诊断（Computer-Aided Diagnosis, CAD）系统提供了新的技术路径。

本课题针对公开的BUSI（Breast Ultrasound Images）数据集——包含780张良性、恶性、正常三类的乳腺超声图像——设计并实现了一套基于深度学习的自动诊断系统。研究采用五阶段递进式实验方法，系统地评估了不同技术方案对分类性能的影响。在基线实验中，基于ResNet18的迁移学习模型在无数据增强条件下达到83.76\%的测试准确率；引入传统在线数据增强后提升至86.32\%；利用DAGAN生成对抗网络合成数据扩充训练集后进一步改善至87.18\%；采用混合增强策略结合Focal Loss损失函数使准确率达到88.89\%，其中正常组织识别率实现从70\%到95\%的质的飞跃；最终方案以EfficientNet-B0为骨干网络，综合运用DAGAN合成数据增强、标签平滑正则化、WeightedRandomSampler均衡采样、$300\times300$高分辨率输入及轻量级增强等优化策略，在117张独立测试集上取得94.87\%的整体准确率和99.63\%的AUC值，各类别准确率分别为良性96.97\%、恶性93.55\%、正常90.00\%。此外，通过Grad-CAM梯度加权类激活映射技术生成了热力图可视化结果，验证了模型的决策依据与医学先验知识一致，确实聚焦于病灶区域进行判断。

实验结果表明，将传统数据增强与GAN合成数据相结合的混合增强策略能够互补性地改善不同类别的学习效果，配合适当的模型架构选型和训练技巧优化，能够在小规模医学图像数据集上训练出兼具高准确性和良好可解释性的深度学习分类模型，为乳腺癌超声图像智能诊断提供了一种可行的技术方案。


\vspace{5pt}

\keywordsch 乳腺癌超声图像；深度学习；卷积神经网络；数据增强；生成对抗网络；计算机辅助诊断
